\section{Introduction}
\label{sec:intro}

Automated guided vehicles (AGVs) are a central part of material handling in smart manufacturing, moving workpieces between storage, processing, and delivery locations. Coordinating a fleet is a joint assignment, routing, and traffic-management problem: vehicles may need the same narrow aisle, workstation approach, or transport resource at overlapping times. A schedule must assign jobs and produce routes and timings that avoid collisions while meeting operational targets such as completion time, delivery due dates, and balanced vehicle utilization.

These requirements become more demanding in multi-floor facilities, where lifts connect otherwise separate floor networks. A lift is a shared resource with finite capacity and non-negligible travel time: a conflict at one lift can delay routes on several floors and propagate through the entire schedule. The prioritized planning baseline of Guo et al. \cite{guo2024} assigns transport tasks and plans routes sequentially by AGV priority, resolving conflicts against previously planned higher-priority vehicles. In that baseline, however, AGVs may wait only at the factory's single parking node. With lifts, the next available floor leg may be blocked by traffic on another floor, and parking-node-only waiting can serialize delays that would be better absorbed near the lift itself.

This paper extends the prioritized-planning baseline with a hybrid architecture for a multi-floor factory. A central server performs job assignment, divides each AGV route into floor-level legs at lift edges, and books lift trips, which remain shared across floors and therefore require system-wide coordination. One edge server per floor plans local route legs using prioritized conflict-free routing, and lift-lobby waiting points let a floor server delay a leg at its entry lobby instead of returning the AGV to parking. Job assignment, lift booking, dispatch, merging, and verification stay centrally orchestrated; the floor servers do not negotiate with each other directly.

The implementation is evaluated against the paper-style centralized configuration (waiting at parking only) and a centralized variant with the same lift lobbies, on ten job sets of 100 jobs and 20 AGVs in a tested three-floor layout with 27 nodes and 4 lifts. Every generated plan passes an independent verifier. The hybrid configuration reduces mean last-delivery makespan by \RedMakespanCAabs\%, total tardiness by \RedTardinessCAabs\%, and load imbalance by \RedImbalanceCAabs\% relative to the paper-style centralized configuration; its busiest-server planning CPU time falls by \RedCpuMaxCAabs\%. These figures describe the tested factory, parameters, and seeds, not general guarantees for all multi-floor factories.

The three narrow contributions are: (i) a corrected reimplementation of the ICARCV 2024 prioritized-planning baseline with independent full-plan verification; (ii) a hybrid central-plus-floor-server architecture that separates central job assignment and lift booking from per-floor route-leg planning; and (iii) lift-lobby waiting points that absorb delays near floor-level conflicts and lift resources.

The remainder is organized as follows. Section~\ref{sec:related} reviews directly related work; Section~\ref{sec:model} defines the scheduling model; Section~\ref{sec:method} describes the algorithms and the hybrid extension; Section~\ref{sec:experiments} reports the setup and results; Section~\ref{sec:conclusion} concludes and states limitations.
